\section{Network Architecture and Associative Retrieval}
\label{sec:architecture}
\subsection{Directed processing pathways}
Figure~\ref{fig:architecture} distinguishes the stored lexical knowledge from its use in neural processing. Its organization follows the distinction between a central knowledge repository and lateral processing pathways illustrated in Figure~1 of \citet{stille2020}, here redrawn for the present lexical core. The central cylinder represents the mental lexicon as the lexical component of world knowledge. It contains the conceptual and phonological graphs, their derived vectors, the lemma identity vectors, and the associations linking the three representations of each lexical entry. ConceptNet and PhonoNet denote the conceptual and phonological graph structures, including their deep-layer components.

\begin{figure}[htbp]
\centering
\includegraphics[width=\linewidth]{figures/architecture.pdf}
\caption{Stored lexical knowledge and its neural realization in the present model. The central cylinder denotes the graphs, vectors, and cross-level associations specifying the mental lexicon within world knowledge. Left: production proceeds from concept through lemma to phonological form. Right: perception proceeds from phonological form through lemma to concept. Each pathway contains separate spiking state buffers and two neural associative memories, one of each variant: WTA selection (competition) and thresholded transfer. Both variants apply thresholding; WTA selection additionally implements competition through lateral inhibition. Corresponding buffers on the two sides use identical stored vectors. Solid arrows indicate neural signal flow; dashed arrows indicate how the shared lexical specification defines the buffers' vocabularies and the memories' associations during network construction. The cylinder is a representation of stored knowledge, not an additional spiking module supplying signals during simulation. WTA: winner-take-all. Original diagram of the lexical processing architecture, informed by \citet[Figure~1]{stille2020}.}
\label{fig:architecture}
\end{figure}

The left and right columns show the neural realization of lexical processing. Production maps a concept to its associated lemma and then to its phonological form; perception maps a phonological form to its lemma and then to its concept. Each column contains three neural state arrays, referred to here as buffers, and two associative memories. The buffers represent the currently active vectors through the collective activity of spiking neurons \citep{eliasmith2013}. The associative memories implement the mappings between levels using the entry correspondences specified in the central repository.

Production and perception are implemented as separate networks, each with its own buffers and associative-memory populations. Corresponding buffers nevertheless use the same vector vocabulary: a given lexical entry has the same stored conceptual vector in both pathways, the same lemma vector in both pathways, and the same phonological form vector in both pathways. These shared vectors are reproduced from the same lexical data and vector seed. The time-dependent neural states are specific to each pathway and its input. Thus, the two columns represent separate neural realizations of a shared lexical specification, rather than two simultaneous directions through one set of buffers.

The dashed links in Figure~\ref{fig:architecture} express this dependence on the lexical specification. During network construction, the stored vectors define the state vocabularies and the input and output vectors of the associative memories. During simulation, retrieval proceeds through the neural connections indicated by solid arrows. Both transformations in each pathway are implemented by heteroassociative memories, which retrieve a stored output vector associated with an input vector from another lexical level. The labels WTA selection (competition) and thresholded transfer in Figure~\ref{fig:architecture} identify two variants of this common associative-memory mechanism. The first memory performs WTA selection (competition) on the way to the lemma buffer; the second performs thresholded transfer from the lemma buffer to the target level. The state self-feedback gain is zero, so the buffers follow their incoming signals, and interlevel processing follows the indicated feedforward direction. Section~\ref{sec:neural-implementation} explains this parameter and its role in the present architecture.

\subsection{Formal specification of associative memories}
Each associative memory links two adjacent lexical levels. Its input vocabulary contains the stored vectors against which the incoming neural state is compared; its output vocabulary contains the vectors that can be retrieved through the corresponding associations. In production, the first memory uses conceptual vectors as its input vocabulary and lemma vectors as its output vocabulary. The second uses lemma vectors as input and phonological form vectors as output. In perception, the first memory maps phonological form vectors to lemma vectors, and the second maps lemma vectors to conceptual vectors. The terms input and output vocabulary therefore refer to the role of a representational level in a particular memory.

For any one of these four mappings, let the columns of $U\in\mathbb R^{D\times N}$ contain the stored input vectors and the columns of $V\in\mathbb R^{D\times N}$ contain their associated output vectors, with $N=613$. Column $i$ in both matrices refers to the same lexical entry. For example, in the concept-to-lemma memory, the column containing the conceptual vector for \emph{Kanne} in $U$ is paired with the column containing its lemma vector in $V$. In the implementation, this correspondence is established through a common entry identifier, or key, in the two vocabularies. NengoSPA's \texttt{by-key} mapping uses these identifiers to assemble the input and output vectors in matching column order. The pairing thus follows lexical identity and remains valid when the vocabularies themselves list their entries in different orders.

Retrieval comprises three operations: comparison with the stored input vectors, selection among the resulting matches, and construction of the associated output. Their functional organization is
\begin{equation}
 \mathbf q(t)=U^{\mathsf T}\mathbf x(t),\qquad
 \mathbf z(t)=\mathcal F[\mathbf q](t),\qquad
 \mathbf y(t)=V\mathbf z(t).
 \label{eq:am}
\end{equation}
Here $\mathbf x(t)$ is the $D$-dimensional vector supplied by the preceding neural buffer. The first operation produces $N$ match values: each $q_i(t)$ is the dot product of the incoming vector with the stored input vector for lexical entry $i$. The neural selection mechanism $\mathcal F$ transforms these matches into the $N$ decoded selection outputs $z_i(t)$. The final operation combines the stored output vectors using these selection outputs as weights, yielding the $D$-dimensional vector $\mathbf y(t)$ for the next buffer. When the selection output for \emph{Kanne} dominates, for example, the retrieved vector is dominated by the representation associated with \emph{Kanne} at the output level.

Figure~\ref{fig:architecture} identifies the two associative-memory variants as \emph{WTA selection (competition)} and \emph{thresholded transfer}. WTA selection (competition) denotes the competitive associative memory used at the first transition in each pathway: concept to lemma in production and phonological form to lemma in perception. Thresholded transfer denotes the associative memory used at the second transition: lemma to phonological form in production and lemma to concept in perception. Both variants use a threshold to admit sufficiently strong matches. The WTA selection variant additionally introduces competition between candidates to select a lemma identity; the thresholded transfer variant passes qualifying lemma matches to their associated target representations. WTA selection performs the cleanup role \citep{stewart2011}: it favors a stored lexical candidate when the input also matches other candidates. Here, cleanup is combined with a mapping to another lexical level, giving a heteroassociative memory. The matrices $U$ and $V$ specify the stored correspondences, and $\mathcal F$ supplies the nonlinear thresholding and selection behavior. Section~4.3 describes the two selection mechanisms in detail.

Equation~\eqref{eq:am} separates these functional operations to make their roles explicit. In the spiking network, signals also undergo temporal filtering at synaptic connections. These connection filters are low-pass filters that shape how neural signals evolve over time \citep{eliasmith2003}; cleanup memories instead operate on the match values to select stored representations. The compact equation leaves the synaptic filter kernels implicit. The resulting time courses depend jointly on these filters, neuronal dynamics, and the selection circuit, with the principal time constants specified in Section~4.4.

\subsection{WTA selection (competition) and thresholded transfer}
The two associative-memory variants shown in Figure~\ref{fig:architecture}, \emph{WTA selection (competition)} and \emph{thresholded transfer}, share the same input-comparison and output-mapping organization. Here, \emph{thresholding} denotes the operation that admits matches above a threshold and is used in both variants. The term \emph{competition} denotes the additional interaction between candidates in WTA selection, implemented through lateral inhibition. Thresholded transfer evaluates each candidate against the threshold independently; WTA selection additionally allows strongly activated candidates to suppress competing responses. In NengoSPA, WTA selection (competition) is implemented by \texttt{WTAAssocMem}, and thresholded transfer by \texttt{ThresholdingAssocMem} \citep{nengospaAM}. These are two variants of associative memory, distinguished by their selection mechanisms.

Both memories use a match threshold of $\theta=0.3$. For the thresholded transfer variant used in the second memory, an idealized static target function is
\begin{equation}
 z_i=\begin{cases}q_i,&q_i>\theta,\\0,&q_i\leq\theta.\end{cases}
 \label{eq:threshold}
\end{equation}
This target function preserves the amplitude of each match that exceeds the threshold. The spiking implementation approximates this behavior with neural dynamics and decoding variability. Several entries can pass when several matches exceed threshold, so this variant supports the transfer of multiple qualifying associations.

The WTA selection (competition) variant used in the first memory combines lateral inhibition with a constant decoded output function of one for a surviving selection. In an ideal settled state, it approximates a single selected output vector rather than preserving the winning input's amplitude. The implementation inhibits each selection population using the other populations' raw thresholded outputs. The inhibition scale is 1.0 and the inhibitory synapse has a 5-ms time constant. Selection transients and ties cannot be described fully by an instantaneous maximum operation.

This arrangement retains the division of roles in the preliminary WWT source implementation developed by Katharina Stille: WTA selection (competition) provides cleanup on the way to the lemma level, followed by thresholded transfer to the target level. The first memory receives a structured conceptual or phonological representation whose shared components can yield matches to several lexical candidates. WTA selection (competition) supports the selection of a distinct lemma identity from these matches. The second memory then retrieves the phonological form or concept associated with that selected lemma. Because lemma vectors encode distinct lexical identities and are approximately orthogonal, thresholded transfer is used at this stage to pass sufficiently strong matches to their associated outputs while preserving their amplitudes.

These distinctions give the activation curves a specific computational interpretation: the WTA selection (competition) memory provides a competitive selection signal, the thresholded transfer memory provides an amplitude-sensitive signal, and the state arrays represent distributed vectors. For example, an entry output near 1 in the WTA selection (competition) memory reflects the configured signal amplitude for a selected candidate. In the thresholded transfer memory, several entries can contribute simultaneously, with each entry's output specifying its weight in the retrieved vector. The curves thus describe selection and signal strength over time. Section~\ref{sec:neural-probabilistic} discusses how this mechanistic account relates to models that explicitly represent uncertainty over lexical alternatives.

\subsection{Neural implementation and parameters}
\label{sec:neural-implementation}
The neural buffers in Figure~\ref{fig:architecture} are implemented as \emph{state arrays}. An \emph{ensemble} is a group of neurons whose collective activity represents a quantity \citep{eliasmith2003}; an \emph{array} is a collection of such ensembles working together. The term \emph{state} refers to the information currently represented by their activity. In the present implementation, each ensemble represents one component of the buffer's current vector, and the ensembles together provide the complete neural representation of the buffer state.

The \emph{state self-feedback gain} specifies the factor by which a buffer's decoded output is scaled when fed back to its own input. Here, the gain is zero: each buffer is driven by the external input or the preceding associative memory. This choice matches the focus on input-driven lexical retrieval. In the perception and production tasks, the input representation is supplied continuously during retrieval and evaluation (Section~\ref{sec:tasks}), sustaining the signals passed through the successive processing stages. The sustained input thus provides the temporal support needed for the present tasks, while the buffers represent and transmit the available signals. After the driving signal subsides, residual activity decays according to the neuronal and synaptic dynamics. The zero gain applies specifically to buffer self-feedback; competition within the WTA associative memories continues to use lateral inhibition.

A possible next extension would give the input buffer a short-term maintenance function through controlled recurrent self-feedback. In perception, a briefly supplied phonological form vector would be retained in the phonological input buffer; in production, the corresponding role would be taken by the conceptual input buffer. The retained activity would continue to drive associative retrieval after the external cue has ended. This would separate the brief presentation of a representation from its internal maintenance and subsequent use, providing a closer match to tasks in which information is available only temporarily. This recurrent input buffer is proposed as a further development; the reported simulations use sustained input during retrieval. Section~\ref{sec:limitations} outlines the corresponding implementation and evaluation steps.

With $D=512$ one-dimensional ensembles and 50 LIF neurons per ensemble, each state array contains $512\times50=25{,}600$ neurons. Each associative memory uses 100 neurons per lexical candidate, giving $613\times100=61{,}300$ neurons. The three buffers and two associative memories in each pathway therefore contain a total of
\begin{equation}
 3(25{,}600)+2(61{,}300)=199{,}400.
\end{equation}
Input-to-state and intermodule connections use low-pass synapses with a 10-ms time constant, which shape the transmission of activity between neural populations. Together with the 5-ms inhibitory time constant of the WTA selection (competition) memory, these synaptic parameters contribute to the dynamics of the neural architecture.

Two separate seeds control the numerical representation and its neural realization. The \emph{vector seed}, fixed at 10, selects the vector vocabulary described in Section~\ref{sec:graph-encoding}. Once that vocabulary has been constructed, the \emph{network seed}, fixed at 11, controls sampled properties of the neural populations, including encoders and neuronal response parameters. Keeping the vectors fixed while changing the network seed produces an alternative neural realization of the same representational geometry and intended associations. Changing the vector seed instead supplies an alternative vocabulary, for which the neural mappings are then constructed. The two seeds therefore make it possible to investigate variation in the numerical encoding separately from variation in the neural implementation. With fixed lexical data, model settings, and software, retaining both seeds reproduces the chosen realization used throughout the reported baseline simulations.

Table~\ref{tab:parameters} summarizes the implemented settings, including the representation sizes, neuron counts, feedback, match threshold, and synaptic time constants. The perception and production tasks applied to this architecture, including trial timing, recording, and evaluation, are specified separately in Section~\ref{sec:tasks}.

\begin{table}[htbp]
\centering\small
\begin{tabular}{ll}\toprule
Quantity & Value\\\midrule
Lexical candidates / vector dimensions & 613 / 512\\
Vector seed / network seed & 10 / 11\\
Dimensions per state ensemble & 1\\
Neurons per state ensemble / memory entry & 50 / 100\\
Neurons per directed pathway & 199,400 LIF\\
State self-feedback gain & 0\\
Associative-memory match threshold & 0.3\\
Intermodule / WTA inhibitory time constant & 10 / 5 ms\\\bottomrule
\end{tabular}
\caption{Representation and neural architecture parameters of the 613-entry model. Library defaults not overridden by the model are preserved by the pinned software environment.}\label{tab:parameters}
\end{table}
\FloatBarrier

